\documentclass[aspectratio=169]{beamer}
\usetheme{Madrid}
\usecolortheme{default}
\usepackage[T1]{fontenc}
\usepackage[utf8]{inputenc}
\usepackage[spanish,es-nodecimaldot]{babel}
\usepackage{amsmath}
\usepackage{booktabs}
\usepackage{graphicx}
\usepackage{siunitx}

\title{Laboratorio 2}
\subtitle{Análisis de ECG y Fibrilación Auricular}
\author{Integrantes: Cristian Stiven Capera C., Ana Sofía García G., Jeimmy Andrea Gonzáles G.,\\
Karol Mariana Gutiérrez G. y Brayan Andrés Ruiz C.}
\date{\today}

\begin{document}

% Diapositiva 1: Portada
\begin{frame}
    \titlepage
\end{frame}

% Diapositiva 2: Problema, usuario, pregunta y límite clínico
\begin{frame}{Problema, Usuario y Límite Clínico}
    \begin{block}{Pregunta de Trabajo}
        ¿Puede un tablero interactivo auditar las discrepancias entre anotaciones clínicas de Fibrilación Auricular (FA) y detecciones algorítmicas bajo condiciones de ruido real?
    \end{block}
    \begin{itemize}
        \small
        \item \textbf{Usuario y Contexto:} Ingenieros biomédicos y clínicos evaluando registros ambulatorios largos del MIT-BIH Atrial Fibrillation Database (AFDB).
        \item \textbf{Problema Técnico:} Los artefactos de movimiento y las fluctuaciones de línea base pueden inducir falsos positivos en detectores automáticos de picos.
        \item \textbf{Límite Clínico:} Los algoritmos basados en umbrales de amplitud no sustituyen la auditoría visual experta ante episodios de mucho ruido.
    \end{itemize}
\end{frame}

% Diapositiva 3: ECG, AFDB y cadena de procesamiento
\begin{frame}{ECG, AFDB y Cadena de Procesamiento}
    \begin{columns}[T]
        \column{0.5\textwidth}
        \textbf{Datos y Arquitectura:}
        \begin{itemize}
            \small
            \item \textbf{AFDB:} Registros de ECG de larga duración con anotaciones de ritmo oficiales en \texttt{data/raw/afdb/}.
            \item \textbf{Procesamiento Dinámico:} Cálculo en memoria RAM sin persistencia pesada en \texttt{data/processed/}.
            \item \textbf{Gestión:} Entorno estricto y determinista administrado mediante \texttt{uv}.
        \end{itemize}
        
        \column{0.5\textwidth}
        \textbf{Filtrado y Trazabilidad:}
        \begin{itemize}
            \small
            \item \textbf{Filtro Pasabanda:} Butterworth de fase cero (0.5--40 Hz) para eliminar la deriva de línea base.
            \item \textbf{Inventario:} Control de integridad bit a bit mediante hashes SHA-256 en archivos JSON de resultados.
        \end{itemize}
    \end{columns}
\end{frame}

% Diapositiva 4: QRS, RR, anotaciones, calidad y pruebas
\begin{frame}{QRS, RR, Anotaciones, Calidad y Pruebas}
    \begin{columns}[T]
        \column{0.5\textwidth}
        \textbf{Detección y Control QRS:}
        \begin{itemize}
            \small
            \item Detección mediante \texttt{xqrs\_detect} sobre la señal procesada.
            \item Control post-detección aplicando restricciones de período refractario y rechazo de duplicados.
            \item Cálculo de intervalos RR excluyendo transiciones de ritmo.
        \end{itemize}
        
        \column{0.5\textwidth}
        \textbf{Garantía y Calidad (\texttt{pytest}):}
        \begin{itemize}
            \small
            \item Validación de contratos de entrada y salida en la configuración de canales.
            \item Cobertura automatizada de pruebas unitarias para métricas y robustez ante ruido.
        \end{itemize}
    \end{columns}
\end{frame}

% Diapositiva 5: Demostración y comparación principal
\begin{frame}{Demostración y Comparación Principal}
    \begin{columns}[T]
        \column{0.48\textwidth}
        \textbf{Exploración Interactiva en Streamlit:}
        \begin{itemize}
            \small
            \item \textbf{Ventanas Cortas:} Permiten visualizar con precisión las cruces de QRS aceptados sobre la señal filtrada.
            \item \textbf{Envolvente Dinámica:} Oculta las marcas QRS en vistas globales amplias para evitar sobrecarga visual y pérdida de rendimiento.
        \end{itemize}
        
        \column{0.48\textwidth}
        \textbf{Evidencia de Limitación en AFDB:}
        \begin{itemize}
            \small
            \item \textbf{Registro 04043 (1100–1130 s):} Bajo episodios de FA anotada con artefactos severos de movimiento, el detector genera marcas falsas sobre oscilaciones espurias.
            \item \textbf{Valor del Dashboard:} Permite auditar visualmente las fallas algorítmicas frente a la referencia clínica oficial.
        \end{itemize}
    \end{columns}
\end{frame}

% Diapositiva 6: Interpretación, limitaciones y reproducción
\begin{frame}[fragile]{Interpretación, Limitaciones y Reproducción}
    \begin{itemize}
        \small
        \item \textbf{Interpretación:} La automatización facilita la revisión masiva, pero la validación humana sigue siendo indispensable ante señales altamente contaminadas.
        \item \textbf{Limitaciones:} Conjunto reducido de registros de la base de datos AFDB y no realiza diagnósticos clínicos en tiempo real, ya que los filtros de fase cero y los algoritmos implementados están diseñados para la auditoría offline.
        \item \textbf{Reproducción Externa:}
    \end{itemize}
    \vspace{-0.1cm}
    \tiny
    \begin{verbatim}
    uv sync --locked && uv run pytest && uv run python scripts/reproduce.py
    latexmk -pdf -interaction=nonstopmode -halt-on-error -outdir=build/slides slides/presentacion.tex
    \end{verbatim}
\end{frame}

\end{document}
